% Zak-Zok — Standard Arabic edition
% STORY POPULATION STARTED from the original scanned book.
% Illustration assignments use the canonical shared-media IDs:
%   1, 2, 3, ... 14
%
% The four page types are modeled explicitly below:
%   TYPE 1 — white page + text only
%   TYPE 2 — full illustration + native LaTeX text
%   TYPE 3 — full illustration, no text
%   TYPE 4 — white page + inset illustration
%
% Build with XeLaTeX.

\documentclass[12pt]{article}

\usepackage{fontspec}
\usepackage{polyglossia}
\usepackage{geometry}
\usepackage{graphicx}
\usepackage{tikz}
\usepackage{ragged2e}
\usepackage{xcolor}

\setdefaultlanguage{arabic}
\setotherlanguage{english}

% Prefer Amiri, with a Noto fallback when available.
\IfFontExistsTF{Amiri}
  {\newfontfamily\arabicfont[Script=Arabic,Scale=1.10]{Amiri}}
  {\newfontfamily\arabicfont[Script=Arabic,Scale=1.10]{Noto Naskh Arabic}}

% ---------------------------------------------------------------------------
% PROVISIONAL BOOK GEOMETRY
% ---------------------------------------------------------------------------
% These dimensions are placeholders for the rebuilt edition.
% Lock the final KDP trim + bleed dimensions before production export.
\newcommand{\BookWidth}{170mm}
\newcommand{\BookHeight}{232mm}

\geometry{
  paperwidth=\BookWidth,
  paperheight=\BookHeight,
  top=14mm,
  bottom=14mm,
  inner=14mm,
  outer=14mm
}

\setlength{\parindent}{0pt}
\setlength{\parskip}{0.55\baselineskip}
\pagestyle{empty}

% ---------------------------------------------------------------------------
% SHARED MEDIA
% ---------------------------------------------------------------------------
% Standard Arabic pages reference approved assets stored in:
%   ../shared-media/
%
% Illustration files are canonically named by sequential ID only:
%   1, 2, 3, ... 14
%
% No extension is hardcoded here. \includegraphics can resolve standard image
% extensions such as .png, .jpg, and .pdf once the assets are uploaded.
\graphicspath{{../shared-media/}}

% ---------------------------------------------------------------------------
% DIALOGUE / DRAMA GEOMETRY
% ---------------------------------------------------------------------------
% The right-hand speaker gutter is reserved.
% Dialogue content begins at a shifted level to its left.
% Wrapped dialogue MUST NOT enter the speaker-name gutter.
% Narration on dialogue-capable pages uses the same body width, so prose also
% cannot slip into the speaker-name spacing level.
\newlength{\SpeakerGutter}
\newlength{\DialogueGap}
\newlength{\StoryBodyWidth}

\setlength{\SpeakerGutter}{28mm}
\setlength{\DialogueGap}{5mm}
\setlength{\StoryBodyWidth}{\dimexpr\textwidth-\SpeakerGutter-\DialogueGap\relax}

% Narration/prose constrained to the body column on dialogue-capable pages.
\newcommand{\Narration}[1]{%
  \par
  \noindent
  \begin{minipage}[t]{\StoryBodyWidth}
    \RaggedLeft
    #1
  \end{minipage}
  \par
}

% Drama-style dialogue:
% body text on the left; fixed speaker gutter on the right.
\newcommand{\DialogueLine}[2]{%
  \par\addvspace{0.75\baselineskip}
  \noindent
  \begin{minipage}[t]{\StoryBodyWidth}
    \RaggedLeft
    #2
  \end{minipage}%
  \hspace{\DialogueGap}%
  \begin{minipage}[t]{\SpeakerGutter}
    \RaggedLeft
    \bfseries #1
  \end{minipage}
  \par
}

% ---------------------------------------------------------------------------
% COMMON IMAGE HELPERS
% ---------------------------------------------------------------------------

% Full-page image layer used by TYPE 2 and TYPE 3.
\newcommand{\FullPageImage}[1]{%
  \begin{tikzpicture}[remember picture,overlay]
    \node[
      anchor=center,
      inner sep=0pt
    ] at (current page.center) {%
      \includegraphics[
        width=\paperwidth,
        height=\paperheight
      ]{#1}%
    };
  \end{tikzpicture}%
}

% ---------------------------------------------------------------------------
% TYPE 1 — WHITE PAGE + TEXT ONLY
% ---------------------------------------------------------------------------
% Use for ordinary text pages.
% #1 = native Arabic LaTeX story content.
\newcommand{\TextOnlyPage}[1]{%
  \clearpage
  \thispagestyle{empty}
  \vspace*{2mm}
  #1
  \clearpage
}

% ---------------------------------------------------------------------------
% TYPE 2 — FULL ILLUSTRATION + TEXT
% ---------------------------------------------------------------------------
% Use when restored artwork fills the page and Arabic is layered above it.
% The Arabic remains native/vector text; it is NOT baked into the image.
%
% #1 = illustration file ID
% #2 = native Arabic LaTeX story content
%
% Page-specific placement can later override the default text block position.
\newcommand{\IllustratedTextPage}[2]{%
  \clearpage
  \thispagestyle{empty}
  \FullPageImage{#1}
  \vspace*{8mm}
  \begin{minipage}[t]{\textwidth}
    #2
  \end{minipage}
  \clearpage
}

% ---------------------------------------------------------------------------
% TYPE 3 — FULL ILLUSTRATION, NO TEXT
% ---------------------------------------------------------------------------
% Use for pure full-page artwork.
% #1 = restored/upscaled illustration file ID.
\newcommand{\FullIllustrationPage}[1]{%
  \clearpage
  \thispagestyle{empty}
  \FullPageImage{#1}
  \null
  \clearpage
}

% ---------------------------------------------------------------------------
% TYPE 4 — WHITE PAGE + INSET ILLUSTRATION
% ---------------------------------------------------------------------------
% Use for a white page carrying a centered inset/vignette illustration.
% #1 = restored/upscaled illustration file ID.
\newcommand{\InsetIllustrationPage}[1]{%
  \clearpage
  \thispagestyle{empty}
  \vspace*{\fill}
  \begin{center}
    \includegraphics[
      width=0.82\textwidth,
      height=0.72\textheight,
      keepaspectratio
    ]{#1}
  \end{center}
  \vspace*{\fill}
  \clearpage
}

\begin{document}

% ===========================================================================
% STORY
% ===========================================================================

% ---------------------------------------------------------------------------
% Original PDF page 9 — printed story page 8 — TYPE 1: white page + text only
% Source transcription preserved from the scanned edition.
% ---------------------------------------------------------------------------
\TextOnlyPage{%
  \begin{center}
    {\Large\bfseries زقزوق و التمساح}
  \end{center}

  \vspace{1.4cm}

  أخذ زقزوق\footnote{زقزوق: أحد الطيور التي توجد بوفرة في الصعيد بين سيول وطيبة واسمه العلمي الزقزاق ذو الرأس السوداء (طير التمساح).} يضرب جناحيه بقوة وكان قلبه الصغير يخفق بشدة وهو طائر فوق النيل العظيم .. كاد في هذه اللحظة أن يبلغ الضفة الشرقية للنيل حيث يسكن هو وإخوته .. كان اليوم هو يوم من أيام الخريف والوقت وقت الغروب والشمس تسحب آخر أشعتها الذهبية من على صفحة النيل، وكان زقزوق ينظر إلى الماء الذي اكتسب لونا أسود غامضا فيه رهبة الليل التي لا ينقذه منها إلا النظر إلى قرص الشمس الأحمر، أو حينما يحول نظره فيرى مدينة «الأقصر» الكائنة في صعيد مصر ومن ورائها قريته الصغيرة التي يحبها.

  أخذ زقزوق يقترب من قريته بسرعة بعدما عبر «الأقصر» شاهد حقولها المربع الجميل وأهلها الطيبين الذين يحبهم. كان كل شيء في القرية يمضي كعادته. لم يتغير شيء في اليومين اللذين ترك فيهما زقزوق قريته.

  ها هو «حمصص» حمار العم «عبد السلام»، عائدا في طريقه من الحقل الصغير حاملا جوالا من القمح.
}

% ---------------------------------------------------------------------------
% Original PDF page 10 — printed story page 9 — TYPE 1: white page + text only
% Source transcription preserved from the scanned edition.
% This page contains dialogue, so narration and dialogue both respect the
% reserved speaker gutter agreed for the rebuilt edition.
% ---------------------------------------------------------------------------
\TextOnlyPage{%
  \Narration{كعادته حياه زقزوق من أعلى نخلة حارة ولم يرد حمصص الذي لم ينتبه، كانت دورته الأخيرة وهو متعب وعائد إلى البيت وخلفه العم «عبد السلام» صاحب الحقل والقرية الصغيرة.}

  \Narration{رأى زقزوق العم «عبد السلام» وتمنى لو نزل من الجو ووقف على كتفيه محييا ومقبلا كما يفعل البني آدمين، ولكن زقزوق حزن في هذه اللحظة حزنا شديدا، لم يحزن لأنه أصغر الكائنات في هذه القرية وإنما حزن لأن العم «عبد السلام» لا يعرفه.}

  \Narration{حيا زقزوق أبا قردان صديق الفلاح. وتمنى أن يكون هو أيضا صديق الفلاح ورد عليه أبو قردان التحية بسرعة وهو واقف على ساق واحدة وكأنه يستعد للنوم بعد يوم شاق من العمل.}

  \Narration{وصل زقزوق إلى آخر الحقل حيث شجرات السرو الثلاث ورأى عشه الصغير على الشجرة الأخيرة، وسمع أصوات إخوته الصغار في العش وهم يتصايحون اطمأن زقزوق عليهم وعلى العش .. لقد انتهى الطعام الذي تركه لهم .. تذكر زقزوق نقص الطعام والغذاء، وجلس بين إخوته وضمهم إلى جناحيه وأغلق جفنيه ونام. إستيقظ زقزوق من فجر اليوم التالي على صوت عصفور صديقه، قام في نشاط وتلفت حوله فوجد العم «عبد السلام» واقفا يصلي. أخذ زقزوق ينحني مثله ويقلده مصدرا صوتا جميلا فيه شكر وفيه دعاء ..}

  \Narration{ترك عصفور شجرته وجاء إلى صديقه زقزوق .. ليتحدث معه كعادتهما.}

  \DialogueLine{عصفور}{زقزوق رجعت أخيرا يازقزوق .. تركتني هنا وحدي يومين.}
  \DialogueLine{زقزوق}{عصفور صديقي .. إني سعيد لرؤيتك.}
  \DialogueLine{عصفور}{احك لى إذا يازقزوق عما رأيته في رحلتك وعما أحضرته لنا.}
  \DialogueLine{زقزوق}{لقد عبرت النيل ياعصفور، وذهبت إلى أسوان وزرت جزيرة النباتات ..}
}

% ===========================================================================
% STORY IMAGE MAPPING — PHOTO NAMES POPULATED, REMAINING STORY TEXT PENDING
% ===========================================================================
%
% These assignments mirror ../shared-media/schema.md.
% The literal image file names are 1, 2, 3, ... 14.
% Calls remain commented until the corresponding image files are uploaded;
% this keeps the backbone compilable before media is present.
%
% ---------------------------------------------------------------------------
% Original page 13 — Photo 1 — TYPE 2: full illustration + text
% ---------------------------------------------------------------------------
% \IllustratedTextPage{1}{%
%   % TODO: populate native Arabic story text for page 13.
% }
%
% ---------------------------------------------------------------------------
% Original page 14 — Photo 2 — TYPE 3: full illustration, no text
% ---------------------------------------------------------------------------
% \FullIllustrationPage{2}
%
% ---------------------------------------------------------------------------
% Original page 17 — Photo 3 — TYPE 2: full illustration + text
% ---------------------------------------------------------------------------
% \IllustratedTextPage{3}{%
%   % TODO: populate native Arabic story text for page 17.
% }
%
% ---------------------------------------------------------------------------
% Original page 18 — Photo 4 — TYPE 3: full illustration, no text
% ---------------------------------------------------------------------------
% \FullIllustrationPage{4}
%
% ---------------------------------------------------------------------------
% Original page 21 — Photo 5 — TYPE 2: full illustration + text
% ---------------------------------------------------------------------------
% \IllustratedTextPage{5}{%
%   % TODO: populate native Arabic story text for page 21.
% }
%
% ---------------------------------------------------------------------------
% Original page 22 — Photo 6 — TYPE 3: full illustration, no text
% ---------------------------------------------------------------------------
% \FullIllustrationPage{6}
%
% ---------------------------------------------------------------------------
% Original page 27 — Photo 7 — TYPE 4: white page + inset illustration
% ---------------------------------------------------------------------------
% \InsetIllustrationPage{7}
%
% ---------------------------------------------------------------------------
% Original page 28 — Photo 8 — TYPE 3: full illustration, no text
% ---------------------------------------------------------------------------
% \FullIllustrationPage{8}
%
% ---------------------------------------------------------------------------
% Original page 29 — Photo 9 — TYPE 4: white page + inset illustration
% ---------------------------------------------------------------------------
% \InsetIllustrationPage{9}
%
% ---------------------------------------------------------------------------
% Original page 33 — Photo 10 — TYPE 2: full illustration + text
% ---------------------------------------------------------------------------
% \IllustratedTextPage{10}{%
%   % TODO: populate native Arabic story text for page 33.
% }
%
% ---------------------------------------------------------------------------
% Original page 34 — Photo 11 — TYPE 3: full illustration, no text
% ---------------------------------------------------------------------------
% \FullIllustrationPage{11}
%
% ---------------------------------------------------------------------------
% Original page 35 — Photo 12 — TYPE 3: full illustration, no text
% ---------------------------------------------------------------------------
% \FullIllustrationPage{12}
%
% ---------------------------------------------------------------------------
% Original page 36 — Photo 13 — TYPE 2: full illustration + text
% ---------------------------------------------------------------------------
% \IllustratedTextPage{13}{%
%   % TODO: populate native Arabic story/dialogue text for page 36.
% }
%
% ---------------------------------------------------------------------------
% Original page 40 — Photo 14 — TYPE 4: white page + inset illustration
% ---------------------------------------------------------------------------
% \InsetIllustrationPage{14}
%
% Other TYPE 1 pages will be populated sequentially from the source scan.

\end{document}
